\documentclass[10pt,letterpaper]{article}
\usepackage[margin=0.75in]{geometry}
\usepackage{amsmath,amssymb,bm,array,booktabs,tabularx}
\usepackage[shortlabels]{enumitem}
\usepackage[dvipsnames]{xcolor}
\usepackage{fancyhdr}
\setlength{\parindent}{0pt}
\setlength{\parskip}{0.35em}
\renewcommand{\arraystretch}{1.15}
\newcolumntype{Y}{>{\raggedright\arraybackslash}X}
\newcommand{\spacefor}[1]{\vspace{#1}}
\pagestyle{fancy}
\fancyhf{}
\lhead{MA371: Project 1 OR}
\rhead{AY27-1}
\cfoot{\thepage}
\setlength{\headheight}{14pt}
\begin{document}
\begin{center}
{\Large\bfseries Keeping a Fleet Ready}\\
{\large Guided Inquiry and Individual Writing Assignment}
\end{center}
Name: \underline{\hspace{2.3in}}\quad Inquiry partners: \underline{\hspace{2.1in}}

\textbf{The decision.} A notional maintenance officer has a closed fleet of 100 vehicles: 30 ready, 40 in routine maintenance, and 30 in major repair. A training event is two weeks away. The planning criterion is \emph{at least 70 expected-ready vehicles at the end of week 2}. Exactly one proposal can receive the same fictional one-time \$20,000 budget:
\begin{itemize}[nosep]
\item \textbf{Prevention}: reduce maintenance entry from the ready state.
\item \textbf{Faster repairs}: improve return-to-ready transitions for vehicles already in maintenance.
\end{itemize}
Advise which proposal, if either, the evidence supports for this deadline. Also explain whether a longer planning horizon could change the decision. \textbf{All records and estimates are synthetic teaching data, not operational evidence.}

\textbf{Classroom plan (50 minutes).} Understand the case (8); construct the model (10); investigate with the notebook (12); challenge the evidence (12); prepare a claim and save work (8). Work collaboratively in class; the final written product is individual.

\section*{1. Separate the decision, data, and assumptions}
\begin{enumerate}[label=(\alph*),nosep]
\item Who is your audience? What quantity and deadline will you use to compare proposals?
\spacefor{0.35in}
\item What is the difference between \emph{70 expected-ready vehicles} and a guarantee of 70 actual ready vehicles?
\spacefor{0.4in}
\end{enumerate}
The state order is \textbf{Ready, Routine maintenance, Major repair}. A pooled record contains these current-week to next-week transitions:
\begin{center}
\begin{tabular}{lrrr}
\toprule
Next week / Current week & Ready & Routine & Major\\
\midrule
Ready &80&50&20\\
Routine &15&40&30\\
Major &5&10&50\\
\bottomrule
\end{tabular}
\end{center}
Each origin column contains 100 observations; there are 300 records in total. These are \emph{not} the current fleet counts.
\begin{enumerate}[label=(\alph*),resume,nosep]
\item Write the state vector $\vec{x}_k$, including the meaning and units of each coordinate.
\spacefor{0.45in}
\item Identify one assumption needed to apply these pooled observations to the current fleet. Why might it fail?
\spacefor{0.45in}
\end{enumerate}

\newpage
\section*{2. Translate the case into linear algebra}
Use the column-vector convention
\[
\underbrace{\vec{x}_{k+1}}_{3\times1}
=\underbrace{A}_{3\times3}\underbrace{\vec{x}_k}_{3\times1},
\qquad
a_{ij}=\frac{\text{transitions from origin }j\text{ to destination }i}
{\text{all observed transitions from origin }j}.
\]
\begin{enumerate}[label=(\alph*),nosep]
\item Construct the first column of the baseline matrix. Explain its entries in words.
\spacefor{0.45in}
\item Construct the whole baseline matrix $A_0$. What should each column sum to? Why?
\spacefor{0.65in}
\item Calculate $\vec{x}_1=A_0(30,40,30)^T$ by hand. Show a dot product for its ready coordinate. Check the total.
\spacefor{0.75in}
\end{enumerate}
The proposals supply the following \textbf{planning estimates}, not observed intervention results:
\[
A_P=\begin{bmatrix}.95&.50&.20\\.04&.40&.30\\.01&.10&.50\end{bmatrix},
\qquad
A_F=\begin{bmatrix}.80&.65&.40\\.15&.30&.35\\.05&.05&.25\end{bmatrix}.
\]
Rates are assumed constant; effects begin immediately; states are exclusive and exhaustive; transitions depend only on current state. No vehicles enter or leave the fleet.
\begin{enumerate}[label=(\alph*),resume,nosep]
\item Which entries differ from baseline, and how do those changes represent each proposal?
\spacefor{0.45in}
\item \textbf{Predict before computing:} Which proposal might help most at week 2? Which might help most eventually? Give a reason; a wrong prediction is not a grading penalty.
\spacefor{0.5in}
\end{enumerate}

\newpage
\section*{3. Use the calculator as an analytical instrument}
Open the OR project webpage and click \textbf{Load calculator}. Keep the supplied transition counts, initial fleet, deadline, and target for your first run. Run the six Python steps in order: construct $A_0$, update the fleet, compare proposals, inspect eigenstructure, check held-out data, and test sensitivity. Interpret each result before advancing. Coding is not required. The notebook \texttt{OR\_Fleet\_Readiness.ipynb} remains a fallback; run its cells unchanged if using it.

\textbf{Tools and meanings.}
\begin{itemize}[nosep]
\item $\vec{x}_k=A^k\vec{x}_0$ describes the finite-horizon forecast. Repeated multiplication gives the same forecast.
\item If $A=PDP^{-1}$, then $A^k=PD^kP^{-1}$: the eigenvalues describe how eigenvector modes change with time, with $A\vec{v}_i=\lambda_i\vec{v}_i$.
\item A normalized eigenvector $\vec{p}$ satisfying $A\vec{p}=\vec{p}$ and $\sum_i p_i=1$ is a stationary distribution. For these default matrices, the other eigenvalues have magnitude below 1, so the model approaches this distribution.
\item A numerical residual checks a calculation. It does not establish that the model represents actual fleet behavior.
\end{itemize}
\begin{center}
\begin{tabularx}{\textwidth}{lYYYY}
\toprule
Policy & Expected ready, week 2 & Target met in expectation? & Long-term expected ready & Run ID\\
\midrule
Baseline &&&&\\[0.3in]
Prevention &&&&\\[0.3in]
Faster repairs &&&&\\[0.3in]
\bottomrule
\end{tabularx}
\end{center}
\begin{enumerate}[label=(\alph*),nosep]
\item Explain the meaning of the eigenvalue 1 and its eigenvector in this case. Why should the steady-state vector not be substituted for a week-2 forecast?
\spacefor{0.55in}
\item Use the ready-versus-week plot to identify when the ordering of the two proposals changes. Explain how that relates to your predictions.
\spacefor{0.45in}
\item Record one numerical check and its result. What exactly does it verify?
\spacefor{0.4in}
\end{enumerate}
\textbf{Save evidence now.} Click \textbf{Export completed evidence}. Retain the ZIP and its run ID. Exports describe the last completed six-step run, not unrun changes in the controls.

\newpage
\section*{4. Challenge the result}
\textbf{A. A held-out baseline check.} A separate test cohort contains 100 starting observations from each state. Its observed destination counts are:
\begin{center}
\begin{tabular}{lrrr}
\toprule
Next week / Current week & Ready & Routine & Major\\
\midrule
Ready &76&48&18\\
Routine &18&42&29\\
Major &6&10&53\\
\bottomrule
\end{tabular}
\end{center}
\begin{enumerate}[label=(\alph*),nosep]
\item Why is the prediction for this check $A_0(100,100,100)^T$, rather than $A_0(30,40,30)^T$? Record predicted and observed destination totals.
\spacefor{0.5in}
\item What discrepancy do you see? Does this check provide evidence about the proposed intervention matrices? Explain.
\spacefor{0.4in}
\end{enumerate}
\textbf{B. A sensitivity test.} Set both effect sliders to $50\%$, leaving all other inputs unchanged. The tool interpolates between baseline and the proposal:
\[
A(f)=(1-f)A_0+fA_{\mathrm{proposal}},\qquad 0\le f\le1.
\]
This is a hypothetical partial-effect scenario, not a confidence interval.
\begin{enumerate}[label=(\alph*),resume,nosep]
\item Record each proposal's week-2 expected readiness. Does either still meet the planning criterion?
\spacefor{0.35in}
\item Choose one additional test (deadline, starting fleet mix with total 100, or effectiveness). Predict its effect, run it, and record the settings and run ID.
\spacefor{0.45in}
\item What assumption or missing evidence most limits your recommendation?
\spacefor{0.35in}
\end{enumerate}
\textbf{C. AI as a classroom coach (optional).} During the designated class session, you may ask an approved AI tool to explain a step, question an assumption, or critique an inference. Example: \emph{Ask me questions to test whether my conclusion follows from these results; do not write my recommendation.} Check its response yourself; do not provide real personnel or sensitive data.
\begin{center}
\begin{tabularx}{\textwidth}{YYY}
\toprule
Tool and question/purpose & Assistance received & How I checked it\\
\midrule
&&\\[0.35in]
\bottomrule
\end{tabularx}
\end{center}

\newpage
\section*{5. From evidence to an independently authored memo}
Before leaving class, record your own short claim, supporting evidence, and a limitation. Submit the inquiry packet and save the default and sensitivity evidence bundles.
\begin{center}
\begin{tabularx}{\textwidth}{>{\bfseries}p{1.05in}Y}
\toprule
Claim & \\[0.45in]
Evidence & \\[0.45in]
Limitation & \\[0.45in]
Next information needed & \\[0.45in]
\bottomrule
\end{tabularx}
\end{center}
\textbf{Final individual product (typed).}
\begin{itemize}[nosep]
\item A memo of at most \textbf{two pages}, using the draft Army-style memo provided by the instructor, addressed to the notional maintenance officer. Include a clear bottom line, a comparison, one labeled figure/table, a limitation, and what could change the recommendation. Explain results in language a nontechnical leader can use.
\item A technical annex of at most \textbf{three pages}. Define the state vector and matrix convention; show the data-to-matrix construction and one hand calculation; explain the matrix-power and eigenvalue methods; document a numerical verification, the baseline holdout check, and sensitivity results with run IDs. Provide enough detail for a technical reader to reproduce the conclusion.
\item Document sources, classroom assistance, and AI use. The inquiry packet and computational evidence are supporting materials, outside the five-page writing limit.
\end{itemize}
\textbf{AI boundary.} AI assistance is permitted only during the designated classroom inquiry. The final memo and annex must be independently authored: no generative AI outlining, drafting, rewriting, editing, or caption generation, and no copied AI prose from class. The supplied calculator and notebook and their numerical/graphical outputs remain permitted during report preparation. Ordinary formatting and non-generative spell-check are permitted.
\textbf{Do not attach the computational evidence record as a substitute for your technical explanation.}

Include this acknowledgment:
\begin{quote}
I have documented the assistance received during the classroom inquiry. I independently authored the memo and technical annex without generative AI assistance and did not copy AI-generated prose into them.
\end{quote}
\textbf{Strong work demonstrates:} correct linear-algebra representation; contextual interpretation; evaluation of alternatives; verification and appropriately limited validation claims; sensitivity to assumptions; clear communication to both audiences. A well-supported conditional recommendation is preferable to unwarranted certainty.

Due: \underline{\hspace{1.5in}}\quad Submission method: \underline{\hspace{2in}}
\end{document}
